\section{Reflections}

\subsection{Technical reflection}
Separating the controller from console input makes rules executable without
simulating the terminal. Timestamp checks are particularly important: the
visually similar READY, SERVED and READY\_TO\_PAY states must still occur on
different turns. The independent rule checks complement complete CLI runs,
which could otherwise repeat an error in both a transcript and its expected output.

The chosen scope keeps the existing small design and deterministic choices.
It does not add a GUI, a generic engine or a new advice system. Those features
are not required to explain the OO design and could increase the verification
burden. Current limitations include concrete dependencies, String-based states
and container copies that still expose domain-object references. These are
explicit trade-offs rather than claims of a fully general production model.

\subsection{Personal reflections}

\subsubsection*{Liang Rui}
My main contribution was the Stage~1 core: the \texttt{GameController} turn loop,
the \texttt{TurnPhase} state constants, and the base domain classes
(\texttt{Customer}, \texttt{Order}, \texttt{Waiter}, \texttt{Chef}, \texttt{Cashier}).

The most interesting design decision was using string constants with a
\texttt{requirePhase()} guard instead of an enum or a switch chain.
This keeps illegal phase transitions a runtime exception rather than a silent
path through the wrong branch.
I revised the original idea of advancing the phase inside each action method ---
those advances turned out to be error-prone when two actions could share the
same guard --- and moved the phase field entirely into \texttt{GameController},
so each public method both checks and sets it in one place.

The hardest part was the payment timing rule: a customer who reaches
\texttt{READY\_TO\_PAY} on turn $N$ cannot pay until turn $N+1$ or later.
Off-by-one errors here caused correct-looking transcripts that failed on the
final turn. I added a dedicated \texttt{readyToPayTurn} field and a
\texttt{canPay(turn)} predicate that uses a strict less-than comparison to
enforce the full-turn gap.

\subsubsection*{Xie Yixuan}
I extended Stage~1 with the \texttt{VIPCustomer} and \texttt{CriticCustomer}
subclasses, the \texttt{ArrivalSchedule} factory, and the Stage~2 class diagram.

The core concept I applied was the template-method pattern for discounts.
Adding a \texttt{discountedBill(int baseCents)} hook to \texttt{Customer} ---
returning the value unchanged by default --- let \texttt{VIPCustomer} override
it with a 10\% reduction without changing any call sites in \texttt{Waiter} or
\texttt{GameController}.
Keeping the override in the subclass, rather than adding a type check in the
caller, is the polymorphism principle from lecture directly at work.

The trickiest part was \texttt{CriticCustomer.chooseMeal()}.
The selection rule is: pick the item with the most preparation units, breaking
ties by highest price.
I initially wrote two separate loops, but that required two passes over the menu.
Combining the condition into a single expression ---
\texttt{units~>~best.units || (units == best.units \&\& price~>~best.price)} ---
let me handle both criteria in one pass and made the intent clearer.

\subsubsection*{Angel Chua}
I built the Stage~3 features: the \texttt{Host} phase, \texttt{ComboMeal},
\texttt{HappyHour}, and \texttt{PricingService}.

The design decision I found most educational was using dependency injection for
\texttt{PricingService}: instead of letting the service create its own
\texttt{HappyHour} instance, it receives one from \texttt{GameController}.
This keeps the service stateless and testable without a running game loop.

Modelling \texttt{ComboMeal} as a subclass of \texttt{MenuItem} rather than a
separate concept also taught me composition within inheritance: the two component
items are stored and exposed, but all callers that work with \texttt{MenuItem}
references require no changes.

The hardest part was the integer-rounding rule.
Using \texttt{(price * 90 + 50) / 100} for a 10\% discount ensures half-cent
values round up consistently.
I initially used a floating-point cast, which produced rounding inconsistencies
against the expected transcript, and switched to the integer formula after
tracing the difference in a failing fixture run.

\subsubsection*{Zhu Yuean}
My contribution spanned all three stages: I wrote the \texttt{.in}/\texttt{.out}
test fixtures, the \texttt{RuleChecks} suite, the sequence diagram, and the
\texttt{results.json} functional evidence.

Writing deterministic test fixtures taught me how much clarity a fixed input
sequence provides.
Each \texttt{.in} file is a complete transcript of player decisions; the matching
\texttt{.out} file is the exact expected output.
Running a stage tag against its fixtures gives an objective pass/fail signal that
does not depend on manual inspection.

Drawing the sequence diagram challenged me to choose the right level of
abstraction.
I settled on the VIP~+~Happy Hour scenario starting at turn~4, because it shows
the full discount chain: \texttt{Waiter.takeOrder()} calls
\texttt{PricingService.quote()}, which calls \texttt{customer.discountedBill()}
and then checks \texttt{happyHour.isActive()}.
Showing both calls in sequence makes the layered discount policy visible in a
way that the class diagram alone cannot.
